% Lecture 12 notes: \input by notes.tex, \subimport'ed by the semester book.
% Exam questions (экзамен): 26, 27, 28 (lattices, Plan 4).
\begin{lecture}{12}{Алгоритм LLL}

\section{Зачем редуцировать базисы решёток}

В лекции~11 мы познакомились с решётками $\Lat = \{a_1 b_1 + \cdots + a_n b_n : a_i \in \ZZ\}$
и двумя основными вычислительными задачами: задачей о кратчайшем векторе (SVP) и задачей
о ближайшем векторе (CVP). Обе задачи легко решаются, если известен \emph{хороший} базис "---
короткий и почти ортогональный, "--- и становятся трудными, если дан только \emph{плохой}
базис из длинных, почти параллельных векторов. В размерности~$2$ хороший базис по плохому
строит алгоритм Гаусса; в этой лекции мы изучим его многомерное обобщение "--- алгоритм
А.~Ленстры, Х.~Ленстры и Л.~Ловаса (LLL, 1982)~\cite{lll}.

Алгоритм LLL работает за \emph{полиномиальное} время и находит базис, первый вектор которого
длиннее кратчайшего вектора решётки не более чем в $2^{(n-1)/2}$ раз. Экспоненциальный
множитель кажется слабым результатом, но для криптоанализа этого часто достаточно:
\begin{itemize}
  \item рюкзачные криптосистемы (Меркле"--~Хеллман и её варианты) и задача о сумме подмножества
    малой плотности решаются одним вызовом LLL;
  \item конгруэнтная криптосистема (лекция~11), система GGH и NTRU малой размерности
    вскрываются поиском короткого вектора в подходящей решётке;
  \item метод Копперсмита находит малые корни многочленов по модулю $N$, что даёт атаки на RSA
    с малой открытой экспонентой и при частично известных простых делителях.
\end{itemize}
С другой стороны, современная постквантовая криптография (стандарты ML-KEM и ML-DSA, 2024)
опирается на то, что в размерностях в несколько сотен ни LLL, ни его улучшения не находят
достаточно коротких векторов.

Всюду ниже векторы "--- строки, $\langle x, y\rangle$ "--- стандартное скалярное
произведение в $\RR^m$, $\|x\| = \sqrt{\langle x, x \rangle}$, $b_1, \dots, b_n \in \RR^m$ "---
линейно независимые векторы, $\Lat = \Lat(b_1, \dots, b_n)$ "--- порождённая ими решётка,
$\lambda_1(\Lat)$ "--- длина кратчайшего ненулевого вектора $\Lat$. Через
$\lfloor x \rceil$ обозначается ближайшее к $x$ целое число.

\section{Ортогонализация Грама"--~Шмидта}

\begin{definition}\label{l12:gs}
\emph{Ортогонализацией Грама"--~Шмидта} базиса $b_1, \dots, b_n$ называются векторы
\[
  b_1^* = b_1, \qquad
  b_i^* = b_i - \sum_{j=1}^{i-1} \mu_{ij}\, b_j^*, \qquad
  \mu_{ij} = \frac{\langle b_i, b_j^* \rangle}{\langle b_j^*, b_j^*\rangle}
  \quad (1 \le j < i \le n).
\]
Обозначим $B_i = \|b_i^*\|^2$; положим также $\mu_{ii} = 1$ и $\mu_{ij} = 0$ при $j > i$.
\end{definition}

В отличие от курса линейной алгебры, векторы $b_i^*$ \emph{не нормируются}: так все величины
$\mu_{ij}$, $B_i$ и координаты $b_i^*$ остаются рациональными для целочисленного базиса.
Векторы $b_i^*$, вообще говоря, не принадлежат решётке.

\begin{lemma}\label{l12:gsprops}
Пусть $V_i = \Span(b_1, \dots, b_i)$, $V_0 = \{0\}$. Тогда
\begin{enumerate}
  \item векторы $b_1^*, \dots, b_n^*$ попарно ортогональны и $V_i = \Span(b_1^*, \dots, b_i^*)$;
  \item $b_i^*$ "--- ортогональная проекция $b_i$ на $V_{i-1}^\perp$; в частности,
    $\|b_i^*\| \le \|b_i\|$ и $b_i^* \ne 0$;
  \item $b_i = \sum_{j \le i} \mu_{ij} b_j^*$, т.\,е. $\mathbf{B} = M \mathbf{B}^*$, где
    $\mathbf{B}$ и $\mathbf{B}^*$ "--- матрицы со строками $b_i$ и $b_i^*$, а
    $M = (\mu_{ij})$ "--- нижняя унитреугольная матрица;
  \item $\mu_{ij} = \langle b_i, b_j^*\rangle / B_j$ "--- координата $b_i$ вдоль $b_j^*$, и
    $\|b_i\|^2 = B_i + \sum_{j<i} \mu_{ij}^2 B_j$.
\end{enumerate}
\end{lemma}

\begin{proof}
(1) Индукция по $i$. Пусть $b_1^*, \dots, b_{i-1}^*$ попарно ортогональны и порождают
$V_{i-1}$. Для $k < i$
\[
  \langle b_i^*, b_k^*\rangle = \langle b_i, b_k^* \rangle - \sum_{j<i} \mu_{ij}
  \langle b_j^*, b_k^*\rangle = \langle b_i, b_k^* \rangle - \mu_{ik} B_k = 0
\]
по определению $\mu_{ik}$. Так как $b_i^* \in b_i + V_{i-1}$, то
$\Span(b_1^*, \dots, b_i^*) = \Span(b_1, \dots, b_{i-1}, b_i) = V_i$.

(2) Мы показали, что $b_i - b_i^* \in V_{i-1}$ и $b_i^* \perp V_{i-1}$; это и есть
определение ортогональной проекции на $V_{i-1}^\perp$. Отсюда
$\|b_i\|^2 = \|b_i^*\|^2 + \|b_i - b_i^*\|^2 \ge \|b_i^*\|^2$; если бы $b_i^* = 0$, то
$b_i \in V_{i-1}$ вопреки линейной независимости.

(3) "--- переписанное определение. (4) Умножим равенство (3) скалярно на $b_j^*$ и воспользуемся
ортогональностью; формула для $\|b_i\|^2$ "--- теорема Пифагора.
\end{proof}

\begin{theorem}\label{l12:det}
$\det \Lat = \prod_{i=1}^n \|b_i^*\|$. В частности, при $m = n$
$\abs{\det \mathbf{B}} = \prod_{i} \|b_i^*\|$, и это произведение не зависит от выбора базиса
решётки.
\end{theorem}

\begin{proof}
Определитель решётки "--- это $\sqrt{\det G}$, где $G = \mathbf{B}\mathbf{B}^T$ "--- матрица
Грама базиса (при $m = n$ $\det G = (\det \mathbf{B})^2$). По лемме~\ref{l12:gsprops}
$\mathbf{B} = M \mathbf{B}^*$ и $\det M = 1$, поэтому
\[
  \det G = \det(M \mathbf{B}^* \mathbf{B}^{*T} M^T) = \det(\mathbf{B}^*\mathbf{B}^{*T})
  = \det \operatorname{diag}(B_1, \dots, B_n) = \prod_i B_i,
\]
так как векторы $b_i^*$ попарно ортогональны. Независимость от базиса: два базиса одной
решётки связаны унимодулярной матрицей $U$ ($\det U = \pm 1$), и
$\det(U\mathbf{B}\mathbf{B}^T U^T) = \det(\mathbf{B}\mathbf{B}^T)$ (лекция~11).
\end{proof}

\begin{corollary}[неравенство Адамара]
$\det \Lat \le \prod_i \|b_i\|$, причём равенство достигается только для ортогонального
базиса.
\end{corollary}

\begin{proof}
$\|b_i^*\| \le \|b_i\|$ по лемме~\ref{l12:gsprops}, и равенство при всех $i$ означает
$b_i = b_i^*$.
\end{proof}

Таким образом, коэффициент Адамара $\HR(\Lat) = (\det\Lat / \prod_i \|b_i\|)^{1/n}$ из
лекции~11 близок к $1$ ровно тогда, когда векторы $b_i$ близки к своим проекциям $b_i^*$.

\begin{example}\label{l12:gsex}
$b_1 = (1, 0, 1)$, $b_2 = (1, 1, 0)$, $b_3 = (0, 2, 1)$. Тогда $b_1^* = b_1$, $B_1 = 2$;
$\mu_{21} = \frac{1}{2}$, $b_2^* = (1, 1, 0) - \frac12 (1, 0, 1) = (\frac12, 1, -\frac12)$,
$B_2 = \frac32$; $\mu_{31} = \frac12$,
$\mu_{32} = \frac{\langle b_3, b_2^* \rangle}{B_2} = \frac{3/2}{3/2} = 1$,
\[
  b_3^* = (0, 2, 1) - \tfrac12 (1, 0, 1) - (\tfrac12, 1, -\tfrac12) = (-1, 1, 1), \qquad B_3 = 3.
\]
Проверка: $B_1 B_2 B_3 = 2 \cdot \frac32 \cdot 3 = 9 = (\det \mathbf{B})^2$, $\det\mathbf{B} = 3$.
\end{example}

\taskexample{lll-gram-schmidt}

Следующая простая оценка "--- ключ к анализу алгоритма LLL: длины проекций $b_i^*$ оценивают
снизу длину \emph{любого} ненулевого вектора решётки.

\begin{lemma}\label{l12:lambda}
Для любого базиса решётки $\Lat$ выполнено $\lambda_1(\Lat) \ge \min_i \|b_i^*\|$.
\end{lemma}

\begin{proof}
Пусть $v = \sum_{i=1}^n a_i b_i \ne 0$, $a_i \in \ZZ$, и $k$ "--- наибольший индекс
с $a_k \ne 0$. По лемме~\ref{l12:gsprops}(3)
$v = a_k b_k^* + \sum_{j<k} c_j b_j^*$ с некоторыми вещественными $c_j$, и в силу
ортогональности $\|v\|^2 \ge a_k^2 B_k \ge B_k \ge \min_i B_i$, так как $\abs{a_k} \ge 1$.
\end{proof}

\checkyourself{Как изменятся векторы $b_i^*$ и числа $B_i$, если переставить векторы базиса
в обратном порядке? Какие из величин $\mu_{ij}$, $B_i$, $\prod_i B_i$ зависят от порядка?}

\section{LLL-приведённые базисы}

Если бы все $\abs{\mu_{ij}}$ были равны нулю, базис был бы ортогональным, и по
лемме~\ref{l12:lambda} кратчайший из векторов $b_i$ был бы кратчайшим вектором решётки.
Добиться $\mu_{ij} = 0$ в решётке, как правило, нельзя, но можно сделать коэффициенты малыми
и потребовать, чтобы длины $B_i$ убывали не слишком быстро.

\begin{definition}\label{l12:reduced}
Пусть $\frac14 < \delta < 1$. Базис $b_1, \dots, b_n$ называется
\emph{LLL-приведённым} (с параметром $\delta$), если
\begin{enumerate}
  \item (\emph{условие размера}) $\abs{\mu_{ij}} \le \frac12$ при всех $1 \le j < i \le n$;
  \item (\emph{условие Ловаса}) $B_k \ge (\delta - \mu_{k,k-1}^2)\, B_{k-1}$ при
    $k = 2, \dots, n$.
\end{enumerate}
Базис, удовлетворяющий первому условию, называется \emph{приведённым по размеру}.
Стандартный выбор (и всюду далее, если не оговорено иное) "--- $\delta = \frac34$.
\end{definition}

Условие Ловаса имеет наглядный смысл. Обозначим через $\pi_k$ ортогональную проекцию на
$V_{k-1}^\perp$, так что $b_k^* = \pi_k(b_k)$. Тогда $\pi_{k-1}(b_k) = b_k^* +
\mu_{k,k-1} b_{k-1}^*$ и условие Ловаса равносильно неравенству
\[
  \|\pi_{k-1}(b_k)\|^2 \ge \delta\, \|\pi_{k-1}(b_{k-1})\|^2 :
\]
перестановка векторов $b_{k-1}$ и $b_k$ не может уменьшить $(k-1)$-й вектор ортогонализации
больше чем в $1/\sqrt{\delta}$ раз. В размерности~$2$ при $\delta = 1$ это в точности условие
$\|b_2\| \ge \|b_1\|$ из алгоритма Гаусса (см.~раздел~\ref{l12:gauss}).

Из условия размера $\mu_{k,k-1}^2 \le \frac14$, поэтому условие Ловаса влечёт
\emph{ослабленное условие}
\begin{equation}\label{l12:weak}
  B_k \ge \bigl(\delta - \tfrac14\bigr) B_{k-1} = \tfrac12 B_{k-1} \quad (\delta = \tfrac34).
\end{equation}

\begin{theorem}\label{l12:props}
Пусть $b_1, \dots, b_n$ "--- LLL-приведённый базис решётки $\Lat$ с $\delta = \frac34$. Тогда
\begin{enumerate}
  \item $B_j \le 2^{i-j} B_i$ при $1 \le j \le i \le n$;
  \item $\|b_j\|^2 \le 2^{i-1} B_i$ при $1 \le j \le i \le n$;
  \item $\|b_1\| \le 2^{(n-1)/2}\, \lambda_1(\Lat)$;
  \item $\|b_1\| \le 2^{(n-1)/4} (\det \Lat)^{1/n}$;
  \item $\prod_{i=1}^n \|b_i\| \le 2^{n(n-1)/4} \det \Lat$, т.\,е. $\HR(\Lat) \ge 2^{-(n-1)/4}$.
\end{enumerate}
\end{theorem}

\begin{proof}
(1) Из~\eqref{l12:weak} $B_{l-1} \le 2 B_l$; применяя это при $l = i, i-1, \dots, j+1$,
получаем $B_j \le 2^{i-j} B_i$.

(2) По лемме~\ref{l12:gsprops}(4), условию размера и пункту~(1)
\[
\begin{aligned}
  \|b_j\|^2 &= B_j + \sum_{l<j} \mu_{jl}^2 B_l \le B_j + \frac14 \sum_{l=1}^{j-1} B_l
  \le \Bigl(2^{i-j} + \frac14 \sum_{l=1}^{j-1} 2^{i-l}\Bigr) B_i\\
  &= \bigl(2^{i-j} + 2^{i-2} - 2^{i-j-1}\bigr) B_i = \bigl(2^{i-j-1} + 2^{i-2}\bigr) B_i.
\end{aligned}
\]
Поскольку $j \ge 1$, $2^{i-j-1} \le 2^{i-2}$, и $\|b_j\|^2 \le 2^{i-1} B_i$.

(3) Из (1) при $j = 1$: $\|b_1\|^2 = B_1 \le 2^{i-1} B_i \le 2^{n-1} B_i$ для всех $i$,
т.\,е. $\|b_1\|^2 \le 2^{n-1} \min_i B_i \le 2^{n-1} \lambda_1(\Lat)^2$ по
лемме~\ref{l12:lambda}.

(4) Перемножая неравенства $B_1 \le 2^{i-1} B_i$, $i = 1, \dots, n$, и пользуясь
теоремой~\ref{l12:det}:
$B_1^n \le 2^{0 + 1 + \cdots + (n-1)} \prod_i B_i = 2^{n(n-1)/2} (\det\Lat)^2$.

(5) Перемножаем неравенства (2) при $j = i$:
\[\prod_i \|b_i\|^2 \le 2^{n(n-1)/2} \prod_i B_i = 2^{n(n-1)/2} (\det \Lat)^2.\]
\end{proof}

\begin{remark}
Для произвольного $\delta \in (\frac14, 1)$ те же рассуждения дают
$\|b_1\| \le \alpha^{(n-1)/2} \lambda_1(\Lat)$ с $\alpha = 1/(\delta - \frac14)$. Аналогично
доказывается, что $\|b_j\| \le 2^{(n-1)/2} \lambda_j(\Lat)$ для всех $j$, где
$\lambda_j$ "--- последовательные минимумы решётки~\cite{lll}.
\end{remark}

Оценка $2^{(n-1)/2}$ достигается лишь на специально построенных решётках. На практике
LLL ведёт себя гораздо лучше: эксперименты Н.~Гама и Ф.~Нгуена~\cite{gama-nguyen} показывают,
что для «случайных» решёток $\|b_1\| \approx 1{,}02^n (\det\Lat)^{1/n}$ вместо гарантированных
$2^{(n-1)/4} \approx 1{,}19^{n}$. Поэтому в размерностях до нескольких десятков LLL часто
находит кратчайший вектор.

\checkyourself{Покажите, что базис $(5, 0)$, $(2, 4)$ LLL-приведён с $\delta = \frac34$, но
$b_1$ не является кратчайшим вектором решётки. Согласуется ли это с
теоремой~\ref{l12:props}(3)?}

\section{Алгоритм LLL}

Алгоритм поддерживает индекс $k$ и инвариант «векторы $b_1, \dots, b_{k-1}$ образуют
LLL-приведённый базис». На каждом шаге вектор $b_k$ сначала приводится по размеру
относительно $b_{k-1}, \dots, b_1$, затем проверяется условие Ловаса: если оно выполнено,
$k$ увеличивается, иначе $b_{k-1}$ и $b_k$ меняются местами и $k$ уменьшается.

\subsection{Уменьшение размера}

\begin{lemma}\label{l12:sizered}
Пусть $j < k$, $q \in \ZZ$, и $b_k$ заменяется на $b_k' = b_k - q b_j$. Тогда
\begin{enumerate}
  \item решётка не меняется, все векторы $b_i^*$ (и числа $B_i$) не меняются;
  \item меняются только коэффициенты $\mu_{k l}$ при $l \le j$:
    $\mu_{kl}' = \mu_{kl} - q \mu_{jl}$ (в частности, $\mu_{kj}' = \mu_{kj} - q$).
\end{enumerate}
\end{lemma}

\begin{proof}
Матрица перехода унимодулярна, поэтому решётка та же. Пространства $V_i$ не меняются
($b_k' \in b_k + V_{k-1}$), поэтому не меняются и проекции $b_i^* = \pi_i(b_i)$ при $i \ne k$, а
$b_k'^* = \pi_k(b_k - q b_j) = \pi_k(b_k) = b_k^*$, так как $b_j \in V_{k-1}$. Наконец,
$\mu_{kl}' = \langle b_k - q b_j, b_l^*\rangle / B_l = \mu_{kl} - q \mu_{jl}$, и
$\mu_{jl} = 0$ при $l > j$, $\mu_{jj} = 1$.
\end{proof}

Значит, выбирая $q = \lfloor \mu_{kj} \rceil$ последовательно для $j = k - 1, k - 2, \dots, 1$,
мы получаем $\abs{\mu_{kj}'} \le \frac12$, и последующие шаги (с меньшими $j$) уже достигнутые
коэффициенты $\mu_{kl}$, $l > j$, не портят. При $\abs{\mu_{kj}} \le \frac12$ полагаем $q = 0$.

\subsection{Обмен}

\begin{lemma}\label{l12:swap}
Пусть векторы $b_{k-1}$ и $b_k$ переставлены, $\mu = \mu_{k,k-1}$, и штрихом обозначены
величины для нового базиса. Тогда $b_i'^* = b_i^*$ при $i \ne k-1, k$,
\[
  b_{k-1}'^* = b_k^* + \mu\, b_{k-1}^*, \qquad B_{k-1}' = B_k + \mu^2 B_{k-1}, \qquad
  B_{k-1}' B_k' = B_{k-1} B_k .
\]
\end{lemma}

\begin{proof}
Пространства $V_i$ при $i \ne k-1$ не меняются (как множества векторов $b_1, \dots, b_i$
с точностью до порядка), поэтому $b_i'^* = \pi_i(b_i) = b_i^*$ при $i \ne k-1, k$. Новый
$(k-1)$-й вектор ортогонализации "--- проекция $b_k$ на $V_{k-2}^\perp$:
$\pi_{k-1}(b_k) = b_k^* + \mu b_{k-1}^*$ (из разложения $b_k = b_k^* + \mu b_{k-1}^* +
\sum_{l < k-1} \mu_{kl} b_l^*$), и $B_{k-1}' = B_k + \mu^2 B_{k-1}$ по ортогональности.
Наконец, по теореме~\ref{l12:det}, применённой к решёткам $\Lat(b_1, \dots, b_k)$ и
$\Lat(b_1, \dots, b_{k-2})$ (обмен их не меняет), $\prod_{i \le k} B_i' = \prod_{i \le k} B_i$
и $\prod_{i \le k - 2} B_i' = \prod_{i \le k-2} B_i$, откуда $B_{k-1}'B_k' = B_{k-1}B_k$.
\end{proof}

Если условие Ловаса нарушено, то $B_{k-1}' = B_k + \mu^2 B_{k-1} < \delta B_{k-1}$: после
обмена $(k-1)$-й вектор ортогонализации становится короче по крайней мере в
$1/\sqrt{\delta}$ раз. На этом и основано доказательство завершаемости.

\subsection{Алгоритм}

\begin{algorithm}[LLL, $\delta = \frac34$]\label{l12:lllalg}
\begin{algorithmic}[1]
\Require базис $b_1, \dots, b_n$ решётки $\Lat \subset \ZZ^m$.
\Ensure LLL-приведённый базис той же решётки.
\State вычислить $b_i^*$, $\mu_{ij}$, $B_i$ (ортогонализация Грама"--~Шмидта)
\State $k \gets 2$
\While{$k \le n$}
  \For{$j = k - 1, k - 2, \dots, 1$} \Comment{уменьшение размера}
    \State $q \gets \lfloor \mu_{kj} \rceil$ \ ($q = 0$, если $\abs{\mu_{kj}} \le \frac12$)
    \If{$q \ne 0$}
      \State $b_k \gets b_k - q b_j$; \ $\mu_{kl} \gets \mu_{kl} - q \mu_{jl}$ при $l < j$;
        \ $\mu_{kj} \gets \mu_{kj} - q$
    \EndIf
  \EndFor
  \If{$B_k \ge (\frac34 - \mu_{k,k-1}^2)\, B_{k-1}$} \Comment{условие Ловаса}
    \State $k \gets k + 1$
  \Else
    \State поменять местами $b_{k-1}$ и $b_k$, пересчитать $b_{k-1}^*, b_k^*$, $B_{k-1}, B_k$
      и затронутые $\mu$
    \State $k \gets \max(k - 1, 2)$
  \EndIf
\EndWhile
\State \Return $b_1, \dots, b_n$
\end{algorithmic}
\end{algorithm}

Ручной счёт обычно ведут так, как в учебнике~\cite[разд.~7.13.2]{hps}: после каждого
изменения базиса пересчитывают ортогонализацию заново. Эффективные реализации (например,
\cite[алг.~2.6.3]{cohen}) обновляют $\mu_{ij}$ и $B_i$ по явным формулам
леммы~\ref{l12:swap} за $\bigO(n)$ операций.

\begin{example}\label{l12:ex2d}
$b_1 = (13, 7)$, $b_2 = (17, 9)$, $\det \Lat = \abs{13 \cdot 9 - 7 \cdot 17} = 2$.
\begin{enumerate}
  \item $k = 2$: $B_1 = 218$, $\mu_{21} = \frac{284}{218} = \frac{142}{109}$, $q = 1$,
    $b_2 \gets (4, 2)$, $\mu_{21} = \frac{33}{109}$. По теореме~\ref{l12:det}
    $B_2 = 2^2 / 218 = \frac{2}{109}$, а
    $(\frac34 - (\frac{33}{109})^2) \cdot 218 \approx 143{,}5$: условие Ловаса нарушено, обмен:
    $b_1 = (4, 2)$, $b_2 = (13, 7)$, $k = 2$.
  \item $k = 2$: $B_1 = 20$, $\mu_{21} = \frac{66}{20} = \frac{33}{10}$, $q = 3$,
    $b_2 \gets (13, 7) - 3(4, 2) = (1, 1)$, $\mu_{21} = \frac{3}{10}$, $B_2 = \frac{4}{20} = \frac15
    < (\frac34 - \frac{9}{100}) \cdot 20 = \frac{66}{5}$: обмен, $b_1 = (1, 1)$, $b_2 = (4, 2)$.
  \item $k = 2$: $B_1 = 2$, $\mu_{21} = 3$, $q = 3$, $b_2 \gets (4, 2) - 3(1, 1) = (1, -1)$,
    $\mu_{21} = 0$, $B_2 = 2 \ge \frac34 \cdot 2$: условие выполнено, $k = 3$, конец.
\end{enumerate}
Ответ: $(1, 1)$, $(1, -1)$ "--- ортогональный базис. Величина $d_1 = B_1$ (см. ниже) менялась
так: $218 \to 20 \to 2$.
\end{example}

\taskexample{lll-trace}

\subsection{Корректность и завершаемость}

\begin{lemma}\label{l12:invariant}
В начале каждой итерации цикла векторы $b_1, \dots, b_{k-1}$ образуют LLL-приведённый
базис порождаемой ими решётки. Если алгоритм завершается, то он возвращает LLL-приведённый
базис решётки $\Lat$.
\end{lemma}

\begin{proof}
Все операции унимодулярны, поэтому решётка не меняется. При $k = 2$ инвариант тривиален
(один вектор). Итерация не меняет $b_1, \dots, b_{k-1}$ при уменьшении размера $b_k$ и
обеспечивает $\abs{\mu_{kj}} \le \frac12$ для всех $j < k$ (лемма~\ref{l12:sizered}). Если
выполнено условие Ловаса, то $b_1, \dots, b_k$ приведены, и инвариант сохраняется для
$k + 1$. Если нет, то после обмена $b_1, \dots, b_{k-2}$ не изменились (и их
ортогонализация тоже), а $k$ уменьшается до $k - 1$ (или остаётся $2$). Алгоритм
останавливается при $k = n + 1$, т.\,е. когда весь базис приведён.
\end{proof}

Для доказательства завершаемости введём \emph{потенциал}. Пусть
\[
  d_i = \prod_{j=1}^{i} B_j = \det\bigl(\langle b_s, b_t \rangle\bigr)_{s,t=1}^{i}
  = (\det \Lat(b_1, \dots, b_i))^2, \qquad D = \prod_{i=1}^{n-1} d_i .
\]
Равенство $\prod_{j \le i} B_j$ и определителя матрицы Грама векторов $b_1, \dots, b_i$
доказано в теореме~\ref{l12:det}. Для целочисленного базиса матрица Грама целочисленна,
поэтому все $d_i$ "--- натуральные числа, и $D \ge 1$.

\begin{theorem}\label{l12:potential}
Уменьшение размера не меняет $D$, а каждый обмен уменьшает $D$ по крайней мере в
$1/\delta = \frac43$ раза.
\end{theorem}

\begin{proof}
Уменьшение размера не меняет чисел $B_i$ (лемма~\ref{l12:sizered}), а значит, и $D$. При
обмене $b_{k-1} \leftrightarrow b_k$ по лемме~\ref{l12:swap} меняются только $B_{k-1}$ и $B_k$,
причём их произведение сохраняется. Поэтому $d_i$ не меняется при $i \le k - 2$ (сомножители те
же) и при $i \ge k$ (произведение $B_{k-1}B_k$ то же), а
\[
  d_{k-1}' = d_{k-2} B_{k-1}' = d_{k-1} \frac{B_{k-1}'}{B_{k-1}}
  = d_{k-1} \frac{B_k + \mu_{k,k-1}^2 B_{k-1}}{B_{k-1}} < \delta\, d_{k-1},
\]
так как обмен происходит только при нарушенном условии Ловаса
$B_k + \mu_{k,k-1}^2 B_{k-1} < \delta B_{k-1}$. Следовательно, $D' < \delta D$.
\end{proof}

\begin{theorem}\label{l12:termination}
Пусть $b_1, \dots, b_n \in \ZZ^m$ и $\|b_i\|^2 \le X$ для всех $i$. Алгоритм LLL с
$\delta = \frac34$ выполняет не более $s = \frac{n(n-1)}{2} \log_{4/3} X$ обменов и не более
$n - 1 + 2s = \bigO(n^2 \log X)$ итераций главного цикла; число арифметических операций над
рациональными числами равно $\bigO(n^3 m \log X)$.
\end{theorem}

\begin{proof}
В начале $d_i = \prod_{j \le i} B_j \le \prod_{j \le i} \|b_j\|^2 \le X^i$, поэтому
$D \le X^{1 + 2 + \cdots + (n-1)} = X^{n(n-1)/2}$. Каждый обмен уменьшает целое
положительное число $D$ в $\frac43$ раза, значит, обменов не больше
$\log_{4/3} X^{n(n-1)/2} = s$. Каждая итерация либо увеличивает $k$ на $1$, либо делает обмен
(и уменьшает $k$ не более чем на $1$); $k$ растёт от $2$ до $n + 1$, поэтому число
увеличений не превосходит $n - 1 + s$, а всего итераций не больше $n - 1 + 2s$. Итерация
стоит $\bigO(nm)$ операций: уменьшение размера "--- до $k - 1$ вычитаний векторов длины $m$
и пересчётов строки $\mu$, обмен с пересчётом по формулам леммы~\ref{l12:swap} "--- $\bigO(n +
m)$ операций.
\end{proof}

Чтобы получить \emph{битовую} полиномиальную сложность, нужно ещё оценить размеры
числителей и знаменателей. Знаменатели всех $\mu_{ij}$ и $B_i$ делят числа $d_j$
(упражнение: $d_{j}\mu_{ij} \in \ZZ$ и $d_{i-1} b_i^* \in \ZZ^m$), так что они ограничены
$X^{n}$; ограниченность числителей и всех промежуточных векторов $b_k$ доказана
в~\cite{lll}. Итог (\emph{без доказательства}): классический алгоритм LLL
выполняет $\bigO(n^5 m \log^3 X)$ битовых операций при школьном умножении. Современные
варианты с арифметикой с плавающей точкой (L$^2$ Нгуена"--~Стеле, реализованный в библиотеке
fplll) работают за $\bigO(n^4 m \log X (n + \log X))$ и \emph{квадратичны} по $\log X$; см.
обзор~\cite[гл.~5]{nguyen-vallee}.

\begin{remark}
При $\delta = 1$ доказательство завершаемости не проходит (потенциал может не убывать
геометрически), и полиномиальность такого варианта неизвестна. На практике берут
$\delta = 0{,}99$: качество результата немного лучше, а время работы остаётся
полиномиальным, хотя и с большей константой ($\log_{1/\delta} X$ обменов).
\end{remark}

\taskexample{lll-noswap}

\checkyourself{Почему после уменьшения размера $b_k$ не меняются числа $B_i$, а после обмена "---
меняются? Какие именно $B_i$ и $\mu_{ij}$ нужно пересчитать после обмена $b_{k-1}
\leftrightarrow b_k$?}

\checkyourself{Почему для целочисленного базиса все $d_i$ "--- натуральные числа? Что
изменится в доказательстве завершаемости для базиса с рациональными координатами?}

\subsection{Связь с алгоритмом Гаусса}\label{l12:gauss}

При $n = 2$ алгоритм LLL устроен так же, как алгоритм Гаусса из лекции~11: приведение
по размеру $b_2 \gets b_2 - \lfloor \mu_{21} \rceil b_1$ "--- это шаг $v_2 \gets v_2 - m v_1$
с $m = \lfloor v_1 \cdot v_2 / \|v_1\|^2 \rceil$, а обмен "--- перестановка $v_1
\leftrightarrow v_2$. Разница только в условии остановки. Так как $b_2 = b_2^* + \mu_{21}
b_1^*$, то $\|b_2\|^2 = B_2 + \mu_{21}^2 B_1$, и условие Ловаса принимает вид
\[
  \|b_2\|^2 \ge \delta\, \|b_1\|^2 .
\]
При $\delta = 1$ это ровно критерий остановки алгоритма Гаусса, который даёт \emph{кратчайший}
вектор решётки. С $\delta = \frac34$ алгоритм LLL может остановиться раньше: базис $(5, 0)$,
$(2, 4)$ LLL-приведён ($\mu_{21} = \frac25$, $\|b_2\|^2 = 20 \ge \frac34 \cdot 25$), но
кратчайший вектор решётки "--- $(2, 4)$, а не $b_1$. Зато, как мы видели, при $\delta < 1$
гарантировано быстрое (геометрическое) убывание потенциала в любой размерности. Сам алгоритм
LLL можно понимать как «алгоритм Гаусса, применяемый к парам соседних векторов
$\pi_{k-1}(b_{k-1})$, $\pi_{k-1}(b_k)$».

\checkyourself{Докажите, что в размерности $2$ LLL-приведённый базис с $\delta = \frac34$
удовлетворяет $\|b_1\|^2 \le \frac43 \lambda_1(\Lat)^2$. (Указание: рассмотрите
$v = a_1 b_1 + a_2 b_2$ с $a_2 \ne 0$ и оцените $\|v\|^2 \ge a_2^2 B_2$.)}

\section{Приближённое решение CVP: алгоритм ближайшей плоскости}

Алгоритм Бабаи из лекции~11 \emph{округлял координаты} целевого вектора в данном базисе.
Для LLL-приведённого базиса Л.~Бабаи~\cite{babai} предложил и проанализировал другой
вариант "--- \emph{алгоритм ближайшей плоскости}. Решётка $\Lat$ лежит в объединении
параллельных гиперплоскостей $H_c = c\, b_n^* + V_{n-1}$, $c \in \ZZ$: для
$x = \sum_i a_i b_i$ имеем $\langle x, b_n^* \rangle / B_n = a_n$. Выбираем ближайшую к $w$
гиперплоскость $H_{c_n}$ и рекурсивно решаем задачу в решётке
$\Lat(b_1, \dots, b_{n-1})$ размерности на единицу меньше.

\begin{algorithm}[Бабаи, ближайшая плоскость]\label{l12:nearestplane}
\begin{algorithmic}[1]
\Require LLL-приведённый базис $b_1, \dots, b_n$ решётки $\Lat \subset \RR^n$, вектор $w \in \RR^n$.
\Ensure вектор $v \in \Lat$, близкий к $w$.
\State $t \gets w$
\For{$i = n, n - 1, \dots, 1$}
  \State $c_i \gets \lfloor \langle t, b_i^* \rangle / B_i \rceil$; \quad $t \gets t - c_i b_i$
\EndFor
\State \Return $v = w - t = \sum_i c_i b_i$
\end{algorithmic}
\end{algorithm}

На шаге $i$ коэффициент разложения $t$ по $b_i^*$ становится по модулю не больше $\frac12$,
и последующие шаги (вычитание $b_j$, $j < i$) его не меняют, так как $b_j \perp b_i^*$
при $j < i$. Поэтому
\begin{equation}\label{l12:residual}
  w - v = \sum_{i=1}^n e_i b_i^*, \qquad \abs{e_i} \le \tfrac12 .
\end{equation}

\begin{theorem}[Бабаи, 1986]\label{l12:babai}
Если базис LLL-приведён с $\delta = \frac34$, то алгоритм~\ref{l12:nearestplane} находит
$v \in \Lat$ с $\|w - v\| \le 2^{n/2} \min_{u \in \Lat} \|w - u\|$.
\end{theorem}

\begin{proof}
Индукция по $n$. Пусть $u \in \Lat$ "--- ближайший к $w$ вектор. При $n = 1$ алгоритм выбирает
ближайшую точку прямой решётки, т.\,е. $v = u$ по расстоянию. Пусть $n \ge 2$.

\emph{Случай 1: $u \notin H_{c_n}$.} Расстояние от $w$ до любой гиперплоскости $H_c$
с $c \ne c_n$ не меньше $\frac12 \|b_n^*\|$ (так как $c_n$ "--- ближайшее целое к
$\langle w, b_n^*\rangle / B_n$), поэтому $\|w - u\|^2 \ge \frac14 B_n$. С другой стороны,
по~\eqref{l12:residual} и теореме~\ref{l12:props}(1)
\[
  \|w - v\|^2 = \sum_{i} e_i^2 B_i \le \frac14 \sum_{i=1}^n 2^{n-i} B_n < 2^n \cdot \frac14 B_n
  \le 2^n \|w - u\|^2 .
\]

\emph{Случай 2: $u \in H_{c_n}$}, т.\,е. $u = c_n b_n + u'$ с $u' \in \Lat' =
\Lat(b_1, \dots, b_{n-1})$. Разложим $w - c_n b_n = e_n b_n^* + w'$, $w' \in V_{n-1}$. Дальше
алгоритм работает с вектором $t = w - c_n b_n$, и так как $\langle t, b_i^* \rangle =
\langle w', b_i^*\rangle$ при $i < n$, он выдаёт $v = c_n b_n + v'$, где $v'$ "--- результат
алгоритма для $w'$ в решётке $\Lat'$ (с LLL-приведённым базисом $b_1, \dots, b_{n-1}$). По
теореме Пифагора и предположению индукции
\[
  \|w - v\|^2 = e_n^2 B_n + \|w' - v'\|^2 \le e_n^2 B_n + 2^{n-1} \|w' - u'\|^2
  \le 2^n \bigl(e_n^2 B_n + \|w' - u'\|^2\bigr) = 2^n \|w - u\|^2 . \qedhere
\]
\end{proof}

\begin{example}
Базис $b_1 = (1, 32)$, $b_2 = (40, 1)$ LLL-приведён: $B_1 = 1025$, $\mu_{21} = \frac{72}{1025}$,
$b_2^* = \frac{1}{1025}(40928, -1279)$, $B_2 = \frac{1279^2}{1025}$. Для $w = (53, 117)$:
$c_2 = \lfloor \langle w, b_2^*\rangle / B_2 \rceil = \lfloor \frac{1579}{1279} \rceil = 1$,
$t = w - b_2 = (13, 116)$; $c_1 = \lfloor \frac{13 + 3712}{1025} \rceil = \lfloor
\frac{149}{41} \rceil = 4$. Итак, $v = b_2 + 4 b_1 = (44, 129)$, $\|w - v\| = \|(9, -12)\| =
15$; это действительно ближайший вектор решётки (проверено перебором в Sage).
\end{example}

\checkyourself{Покажите, что для ортогонального базиса алгоритм ближайшей плоскости и
округление координат по Бабаи (лекция~11) дают один и тот же, причём ближайший, вектор.}

\section{Применения в криптоанализе}

\subsection{Задача о сумме подмножества и рюкзачные криптосистемы}

\emph{Задача о сумме подмножества}: даны натуральные $M_1, \dots, M_n$ и $S$; найти
$x \in \{0, 1\}^n$ с $\sum_i x_i M_i = S$. В общем случае она NP-трудна, и на ней построены
рюкзачные криптосистемы (Меркле"--~Хеллман, 1978): секретная «лёгкая» (сверхрастущая)
последовательность маскируется умножением по модулю, открытый ключ "--- $M_1, \dots, M_n$,
шифртекст "--- $S$~\cite[разд.~7.2]{hps}. Рассмотрим решётку со строками
\[
  \mathbf{B} = \begin{pmatrix}
    2 & 0 & \cdots & 0 & M_1 \\
    0 & 2 & \cdots & 0 & M_2 \\
    \vdots & & \ddots & & \vdots \\
    0 & 0 & \cdots & 2 & M_n \\
    1 & 1 & \cdots & 1 & S
  \end{pmatrix}.
\]
Если $x$ "--- решение, то $\sum_i x_i b_i - b_{n+1} = (2x_1 - 1, \dots, 2x_n - 1, 0)$ "---
вектор решётки длины $\sqrt n$, тогда как $\det \Lat \approx 2^n S$, и типичная длина
кратчайшего вектора (эвристика Гаусса, лекция~11) значительно больше $\sqrt n$, если числа $M_i$
велики. \emph{Плотность} задачи $d = n / \log_2 \max M_i$. Дж.~Лагариас и Э.~Одлыжко
доказали~\cite{lagarias-odlyzko}, что при $d < 0{,}6463$ почти все задачи сводятся к SVP
в такой решётке; последующие улучшения дали порог $0{,}9408$. На практике для малых $n$
и малой плотности вектор находит уже LLL; для больших "--- BKZ.

\begin{example}
$M = (612354, 823741, 504219, 931577, 718260, 555013)$, $S = 2603163$, $n = 6$, плотность
$d \approx 0{,}30$. LLL ($\delta = \frac34$) для решётки размерности $7$ выдаёт первым вектором
$(-1, 1, -1, -1, 1, -1, 0)$, откуда $x = (1, 0, 1, 1, 0, 1)$; действительно,
$612354 + 504219 + 931577 + 555013 = 2603163$ (проверено в Sage; трасса содержит более
ста итераций "--- это работа для компьютера).
\end{example}

\subsection{Конгруэнтная криптосистема, GGH и NTRU}

В лекции~11 закрытый ключ $(f, g)$ конгруэнтной криптосистемы находился как короткий
вектор решётки с базисом $(1, h)$, $(0, q)$ "--- это LLL (или Гаусс) в размерности $2$.
Так же устроены атаки на системы, где закрытый ключ "--- хороший базис, а открытый "---
плохой базис той же решётки:
\begin{itemize}
  \item \emph{GGH} (Голдрейх"--~Гольдвассер"--~Халеви, 1997): расшифрование "--- алгоритм
    Бабаи с хорошим базисом. LLL и BKZ, применённые к открытому базису (с дополнительными
    структурными наблюдениями), позволили Ф.~Нгуену вскрыть предложенные авторами
    тестовые шифртексты в размерностях до $350$~\cite{nguyen-ggh}.
  \item \emph{NTRU}: закрытый ключ $(f, g)$ "--- короткий вектор решётки размерности $2N$
    с базисом $\begin{pmatrix} I & H \\ 0 & qI \end{pmatrix}$, где $H$ "--- циркулянт открытого
    ключа~\cite[разд.~7.11]{hps}. Для учебных $N$ LLL находит $(f, g)$; для используемых на практике
    $N$ (несколько сотен) ни LLL, ни BKZ не справляются.
\end{itemize}

\subsection{Метод Копперсмита}

\begin{theorem}[Копперсмит, 1996; \emph{без доказательства}]\label{l12:coppersmith}
Пусть $N$ "--- натуральное число, $f(x) \in \ZZ[x]$ "--- унитарный многочлен степени $d$.
Все целые $x_0$ с $f(x_0) \equiv 0 \pmod N$ и $\abs{x_0} < N^{1/d}$ можно найти за время,
полиномиальное по $d$ и $\log N$.
\end{theorem}

Идея доказательства~\cite{coppersmith}: из многочленов $x^j N^{k-i} f(x)^i$, имеющих корень
$x_0$ по модулю $N^k$, строится решётка; LLL находит в ней короткий вектор "--- многочлен
$g(x)$ с малыми коэффициентами, для которого $\abs{g(x_0)} < N^k$, т.\,е. $g(x_0) = 0$ уже
\emph{над $\ZZ$}, и корень находится обычными методами. Следствия: при $e = 3$ сообщение RSA,
у которого известно всё, кроме $1/3$ бит («стереотипное сообщение»), восстанавливается по
шифртексту; модуль $N = pq$ раскладывается, если известна половина старших бит $p$.

\section{Улучшения LLL и практическое поведение}

\begin{itemize}
  \item \emph{BKZ} (Шнорр"--~Ойхнер~\cite{schnorr-euchner}) с размером блока $\beta$:
    вместо пар соседних векторов в блоках из $\beta$ векторов точно решается SVP
    (перебором или просеиванием). При $\beta = 2$ это LLL, при $\beta = n$ "--- точное SVP.
    Качество растёт с $\beta$, а время "--- как $2^{\Theta(\beta)}$ (просеивание) или
    $\beta^{\Theta(\beta)}$ (перебор).
  \item \emph{Практика}~\cite{gama-nguyen}: \emph{корневой коэффициент Эрмита}
    $(\|b_1\| / \det\Lat^{1/n})^{1/n}$ для LLL близок к $1{,}0219$, для BKZ-20 "--- к $1{,}0128$
    и почти не зависит от решётки. Это позволяет предсказывать, какие параметры
    криптосистемы взламываются.
  \item \emph{Квантовые компьютеры}. Для задач SVP/CVP не известно квантовых алгоритмов,
    существенно лучших классических: лучшие квантовые варианты просеивания ускоряют лишь
    константу в показателе ($2^{0{,}265\beta}$ против $2^{0{,}292\beta}$). Поэтому
    решёточные схемы (ML-KEM на основе Kyber, ML-DSA на основе Dilithium, стандарты NIST 2024)
    считаются постквантовыми; их параметры выбирают так, чтобы атака требовала BKZ
    с $\beta \gtrsim 400$.
\end{itemize}

\checkyourself{Почему от криптосистемы на решётках требуется, чтобы размерность решётки
была большой (сотни), а от RSA "--- чтобы был большим модуль? Что именно растёт
экспоненциально в каждом случае?}

\checkyourselfsection

\begin{furtherreading}
  \readitem[разд.~7.13--7.14]{hps}{ортогонализация Грама"--~Шмидта, LLL, его свойства и
    применения к конгруэнтной системе, рюкзакам, GGH и NTRU; разд.~7.2 "--- рюкзачные
    криптосистемы}
  \readitem[гл.~17, разд.~18.1--18.2, гл.~19]{galbraith}{LLL с полным анализом сложности;
    алгоритм ближайшей плоскости; метод Копперсмита}
  \readitem[гл.~2]{micciancio-goldwasser}{приближённые алгоритмы для SVP и CVP: LLL и
    алгоритм ближайшей плоскости}
  \readitem[разд.~2.6]{cohen}{LLL в целочисленной арифметике и формулы пересчёта}
  \readitem[гл.~1, 2, 5, 10]{nguyen-vallee}{история LLL, постоянная Эрмита и анализ LLL,
    варианты с плавающей точкой (L$^2$), применения к RSA (метод Копперсмита)}
  \readitem{lll}{оригинальная статья}
  \readitem{babai}{алгоритм ближайшей плоскости и оценка $2^{n/2}$}
  \readitem{lagarias-odlyzko}{атака на задачи о сумме подмножества малой плотности}
  \readitem{gama-nguyen}{экспериментальное поведение LLL и BKZ}
  \readitem[лекция~2]{regev-lattices}{конспект О.~Регева: алгоритм LLL}
\end{furtherreading}
\end{lecture}
